% Draft discussion.
\paragraph{Strength of the evidence.}
Table~\ref{tab:evidence} states for each main statement whether it was pre-specified, exploratory with a pre-stated rule, or descriptive, which uncertainty is reported, and the main caveat. Only the first row stems from the pre-specified analysis. All other rows are exploratory, and agreement among them is not independent confirmation, because later rules were written after earlier results were known.

\begin{table}[t]
\centering
\small
\begin{tabular}{@{}p{0.19\textwidth}p{0.15\textwidth}p{0.26\textwidth}p{0.28\textwidth}@{}}
\toprule
Statement & Status & Uncertainty reported & Main caveat \\
\midrule
Untrained below Backprop under ridge at V1 (Section~\ref{sec:signflip}) & Pre-specified; verdict ``effect disappears'' & 5 seeds; one-sided sign-flip test; $\Delta_{\min}$; no interval over stimuli, voxels or subjects & One layer; sub-03 near zero; not blind \\
Raw RSA tracks gain coupling (Section~\ref{sec:gainaxis}) & Exploratory; G1 and GEN1 met & Spearman over 70 and 64 cells, no interval; cells not independent & Correlational; G2 not met \\
V1 ranking not identifiable (Section~\ref{sec:multiverse}) & Exploratory; GEN2 met; cell counts descriptive & Stimulus-bootstrap 95\% intervals per cell, uncorrected; cells dependent & 36 chosen cells, not exhaustive; three subjects \\
Angular component favours trained networks (Section~\ref{sec:gaingeometry}) & Exploratory; SEP3 and SEP4 met, SEP2 not met & Stimulus-bootstrap 95\% intervals; counts of models descriptive & ``Geometry'' is a definition; one weight set per ImageNet model; convnext\_tiny exception \\
Resolution slope of own CNN vanishes under centring or standardisation (Sections~\ref{sec:standardization}, \ref{sec:multiverse}) & Exploratory; R3, O2 and P\_ridge met & Stimulus-bootstrap intervals; seed SD for ridge & Rests on three rules together; ImageNet slope remains (GEN4 not met) \\
Two-feature baseline above every network at V1 (Section~\ref{sec:signflip}) & Descriptive & Point estimates, mean over three subjects & No interval for the baseline; neural gain and BOLD not separable \\
\bottomrule
\end{tabular}
\caption{Strength of the evidence for the main statements. ``Exploratory'' denotes analyses with rules committed before their step was run, written after earlier results were known.}
\label{tab:evidence}
\end{table}

\paragraph{What the V1 advantage measured.}
The advantage of the untrained network at V1 that we reported earlier \citep{evalres} was not robust to the analysis choices examined here. Under correlation distance it disappeared when model units were standardised over the stimulus set, and a pre-specified ridge readout placed the untrained network below the trained one. Under Euclidean distance it persisted, with a resolution slope that was not resolved, but the angular component of the Euclidean distance favoured the trained networks. It was also not shared by the untrained network with calibrated batch-normalisation statistics. Raw correlation-distance RSA, the
convention of the earlier work, tracked how strongly a model's per-stimulus gain axis
covaried with that of the voxels, both across our networks and across six ImageNet
architectures. The earlier result is therefore not a statement about representational
geometry; under raw correlation distance it covaries closely with gain coupling. We
could not, however, attribute the differences between rankings to gain alone: removing
the gain axis by rank regression left part of the resolution slope (G2), and the norm
component did not favour the untrained network with a resolved interval in our own
network (SEP2). Under our pre-stated rule, the overall V1 ranking remains unattributed.

\paragraph{A two-feature baseline exceeds every network at V1 in these data.}
Under the ridge readout, a bank of two features, mean luminance and contrast, predicted
V1 better than every network condition at every pooling grid, including the
backpropagation-trained network. In this data set, and for the Conv1 features tested here, the variance in V1
responses that a linear readout captures is thus largely shared with low-level stimulus
statistics, and no network captured it better than the statistics themselves. This comparison rests on point estimates (mean over three subjects, five seeds per network condition) without an interval for the baseline; per-subject values are in Appendix~\ref{app:persubject}. This
limits what V1 comparisons in these data can reveal about learning rules. This is consistent with the gain account: the voxel gain axis
correlated with luminance and contrast, and the untrained network's gain axis shared
these signs. Whether this gain reflects neural mechanisms such as contrast normalisation
\citep{carandini2012} or properties of the BOLD measurement cannot be decided with
these data; the norm of the voxel response vectors did not depend on luminance or
contrast. Data with controlled contrast manipulations or electrophysiological recordings
would be needed to separate the two.

\paragraph{Evaluation resolution.}
In our CIFAR-trained networks, the dependence of the V1 gap on evaluation resolution
vanished under unit centring and standardisation and was absent in the ridge readout.
In the ImageNet models, a small dependence on resolution above the training resolution
remained under standardisation and resided in the angular component. Evaluation
resolution can therefore affect the angular component of model--brain similarity, but
in our earlier work the size of the effect depended on uncentred units. The two model families differ in the ratio
of evaluation to training resolution (up to sevenfold versus twofold), and we do not
claim that the results transfer between these regimes.

\paragraph{Recommendations.}
Five practices follow from these results. (1) Centre units over the evaluated stimulus
set on both the model and the brain side before computing correlation or cosine
distances; normalising voxels over a larger set of trials is not equivalent. (2) Report
gain and geometry components separately, for example with the norm and angle
decomposition used here. (3) Report a small multiverse of distance measure and
normalisation for any ranking claim. (4) Include trivial stimulus-statistic baselines
alongside network models. (5) State how the normalisation statistics of untrained
networks were set, since calibrated and uncalibrated variants behaved differently.

\paragraph{Consequences for our earlier work.}
The headline of \citet{evalres}, that untrained networks exceed backpropagation at V1
at 224\,px and that this gap depends on evaluation resolution, and the corresponding
claim in \citet{leutenegger2026} are not robust to the analysis choices examined here. An audit
of our earlier preprints identified the affected statements; revised versions refer to
this work.
% TODO: confirm once the revised versions are posted; list them if needed.

\paragraph{Limitations.}
The data comprise three subjects from a single data set, and the intervals reported here describe different sources of variation (Section~\ref{sec:inference}): stimulus resampling for RSA and training seeds for ridge. Neither includes variation across subjects, and for each ImageNet model we used one set of trained weights. Ridge scores in one subject were close to zero in every condition, and the primary ridge result is a mean over all three subjects (Appendix~\ref{app:persubject}). The primary analysis was not blind, several pre-specified instruments
were not implemented, and the exploratory analyses were developed sequentially, with
intervals not corrected for multiple comparisons. Five pre-stated predictions failed
(Appendix~\ref{app:failed}); among them, the LOC ranking was not resolved under
Euclidean distance in the ImageNet models (GEN3). One ImageNet model, convnext\_tiny,
showed no V1 effect in either component; its stem normalises activations across
channels at each position \citep{liu2022}, which may remove the offset and norm differences studied
here, but we did not test this. Our use of ``geometry'' for the angular component is a
definition, not a claim about neural coding. Layers were fixed in advance and not
selected for fit; other layers may behave differently.

\paragraph{Conclusion.}
The apparent superiority of untrained networks at V1 in our earlier work is not a
closer match in representational geometry; under the convention we used, it tracks the
coupling of gain axes between model and brain. Centring units over the evaluated
stimuli, separating gain from geometry and reporting trivial baselines make such
comparisons interpretable.
